Write \(r=1/4\).  Consider a design whose early and late binary signals are \(x\) and \(y\), respectively.  Conditional independence of the proxies given \(U\) and Bayes' rule give
\[
 \det(T_d)=\frac{xy(1-r^2)}{1-r^2x^2}>0. \tag{1}
\]
The early- and late-proxy completeness matrices have determinants proportional to \(x\) and \(y\), so they too are nonsingular.  Direct conditional expectation verifies the displayed bridges
\[
 h_d(l,a)=\frac l y,
 \qquad
 q_d(e,a)=\frac{1-r(2a-1)e/x}{1-r^2}. \tag{2}
\]
For example,
\[
 \mathbb E_P[U\mid W_y=l,T=t,G=O]
 =\frac{rt+yl}{1+rtyl},
\]
and substituting this identity in the conditional expectation of \(q_d(W_x,a)\) given \((W_y,T,G=O)\) gives \((1+rtyl)^{-1}\), which is the experimental-to-observational density ratio.  Thus selection-bridge existence and operator bijectivity hold.  Every latent cell has positive probability, so both treatment-overlap conditions and source support hold.  Consistency, latent exchangeability, experimental randomization, full-short-term external validity, and the sequential restriction follow from the construction and the stated conditional independences.  Finite support makes both bridges bounded for every fixed \(\epsilon>0\) and gives moments of every order.  Hence \(P_{\epsilon}\in\mathfrak P_{\mathcal D}\).

Because \(Y=U\), least-squares projection on the two measured proxies gives
\[
 R(x,y)=\mathbb E_P[\operatorname{Var}(U\mid W_x,W_y)\mid G=O]
 =\frac{(1-x^2)(1-y^2)}{1-x^2y^2}. \tag{3}
\]
Consequently
\[
 R_{d_{\mathrm e}}=\frac35,
 \qquad
 R_{d_{\mathrm p}}=\frac{7(1-\epsilon^2)}{16-9\epsilon^2},
\]
and
\[
 R_{d_{\mathrm e}}-R_{d_{\mathrm p}}
 =\frac{13+8\epsilon^2}{5(16-9\epsilon^2)}
 \ge\frac{13}{80}. \tag{4}
\]
Both designs measure two unit-cost variables, so their experimental costs are \(3\), their observational costs are \(4\), and their comparator sample sizes are \(B/7\).  Equation (4) therefore gives
\[
 \Psi_{\mathrm{CLW},d_{\mathrm p}}(P_{\epsilon};B)
 =\frac{7(1-\epsilon^2)(7/16)}{B(1-9\epsilon^2/16)}
 <\frac{21}{5B}
 =\Psi_{\mathrm{CLW},d_{\mathrm e}}(P_{\epsilon};B),
\]
so \(\mathcal D_{\mathrm{pred}}(P_{\epsilon})=\{d_{\mathrm p}\}\).

We next compute the gradient variances.  Since \(p_{g,a}=1/2\), (2) and conditional independence give
\[
 V_E(x,y)=\frac4{y^2}, \tag{5}
\]
and, because \(\mathbb E_P[(U-W_y/y)^2]=(1-y^2)/y^2\),
\[
 V_O(x,y)=\frac4{(1-r^2)^2}\frac{1-y^2}{y^2}
 \left(1-2r^2+\frac{r^2}{x^2}\right). \tag{6}
\]
For \(d_{\mathrm p}\), equations (5)--(6) yield
\[
 V_{O,d_{\mathrm p}}
 =\frac4{(15/16)^2}\frac79
 \left(\frac78+\frac1{16\epsilon^2}\right)
 =\frac{4(1/16)(7/9)}{(15/16)^2\epsilon^2}+O(1), \tag{7}
\]
which implies the stated, weaker remainder \(O(\epsilon^{-1})\).  For \(d_{\mathrm e}\),
\[
 V_{E,d_{\mathrm e}}=16,\qquad
 V_{O,d_{\mathrm e}}=\frac{384}{25},\qquad
 \mathcal V_{d_{\mathrm e}}^*(P_{\epsilon})
 =\left(4\sqrt3+\frac{16\sqrt6}{5}\right)^2<225. \tag{8}
\]
On the other hand,
\[
 \mathcal V_{d_{\mathrm p}}^*(P_{\epsilon})
 \ge4V_{O,d_{\mathrm p}}
 \ge\frac{1792}{2025\epsilon^2}. \tag{9}
\]
For \(0<\epsilon<10^{-3}\), (8)--(9) make \(d_{\mathrm e}\) uniquely efficiency-optimal.  They also show
\[
 \rho_{\mathrm{pred}}(P_{\epsilon})
 =\frac{\mathcal V_{d_{\mathrm p}}^*(P_{\epsilon})}
 {\mathcal V_{d_{\mathrm e}}^*(P_{\epsilon})}
 \longrightarrow\infty. \tag{10}
\]

It remains to verify the neighborhood assertion.  Let
\(\delta=\|P-P_{\epsilon}\|_{\mathrm{cell}}\).  On the fixed witness support, every conditional probability, bridge coefficient, prediction risk, and gradient variance is obtained from the complete cell vector using finite sums, divisions by positive marginal probabilities, and inversion of a \(2\)-by-\(2\) conditional-expectation matrix.  Along the witness path every complete positive cell is at least \(3\times10^{-4}\).  The determinant in (1) is bounded below by \(0.7\epsilon\) for \(d_{\mathrm p}\), while all determinants for \(d_{\mathrm e}\) and all prediction-projection denominators are bounded below by \(0.2\).  Expanding the finite sums, using
\[
 \left|\frac ab-\frac{a'}{b'}\right|
 \le \frac{|a-a'|}{|b|}+\frac{|a'|\,|b-b'|}{|b|\,|b'|}
\]
and the explicit \(2\)-by-\(2\) inverse formula, gives the following conservative bounds whenever \(\delta<10^{-6}\epsilon^4\):
\[
\begin{aligned}
 \max_{d\in\{d_{\mathrm p},d_{\mathrm e}\}}
 |R_d(P)-R_d(P_{\epsilon})|&\le10^9\delta,\\
 \max_{g\in\{E,O\}}|V_{g,d_{\mathrm e}}(P)-V_{g,d_{\mathrm e}}(P_{\epsilon})|
 &\le10^{12}\delta,\\
 \max_{g\in\{E,O\}}|V_{g,d_{\mathrm p}}(P)-V_{g,d_{\mathrm p}}(P_{\epsilon})|
 &\le10^{12}\delta\epsilon^{-3}.
\end{aligned} \tag{11}
\]
These inequalities are simply finite-dimensional perturbation bounds: a conditional-probability entry changes by order \(\delta\), a design-
\(d_{\mathrm p}\) inverse changes by order \(\delta\epsilon^{-2}\), and its quadratic score moment changes by order \(\delta\epsilon^{-3}\).  The displayed numerical constants dominate the bounds obtained from the cell lower bounds above.

Since \(\epsilon<10^{-3}\), (11) perturbs either prediction risk by less than \(10^{-9}\), either variance for \(d_{\mathrm e}\) by less than \(10^{-6}\), and either variance for \(d_{\mathrm p}\) by at most \(10^6\epsilon\le10^3\).  Thus the fixed predictive margin \(13/80\) in (4) remains strict.  Also, (7)--(9) leave \(\mathcal V_{d_{\mathrm p}}^*(P)>8\times10^5\), whereas continuity of the profiled formula and (8), (11) give \(\mathcal V_{d_{\mathrm e}}^*(P)<226\).  Hence throughout the ball the predictive optimizer is \(d_{\mathrm p}\), the efficiency optimizer is \(d_{\mathrm e}\), and the two optimizer sets are disjoint.  Membership in \(\mathfrak P_{\mathcal D}\) is part of the definition of \(\mathcal U_{\mathrm{bin}}\), so every law in the region is regular for both designs.  Finally, every path law \(P_{\epsilon}\) itself belongs to its corresponding ball because its cell distance is zero.  Equation (10) therefore proves
\[
 \sup_{P\in\mathcal U_{\mathrm{bin}}}\rho_{\mathrm{pred}}(P)=\infty.
\]